\documentclass[fontsize=12pt, master]{diploma}

\usepackage{task}

\student{Бирюков Виктор Владимирович}
\faculty{№ 8 <<Компьютерные науки и прикладная математика>>}
\department{\textbf{Образовательный центр Института № 8}}
\speciality{01.04.02 Прикладная математика и информатика}
\profile{Продуктовая разработка и разработка программных систем}
\group{М8О-208СВ-24}

\theme{Предотвращение взаимных блокировок автономных агентов методами централизованного планирования}
\workDate{\uline{\hspace{1cm}}\ \uline{\hspace{2.5cm}}\ \the\year\ г.}
\similarityPercentage{2.05}

\supervisor{Судаков Владимир Анатольевич}
\supervisorCredentials{доктор технических наук, профессор, профессор образовательного центра Института №8 <<Компьютерные науки и прикладная математика>>}

\reviewer{Осипов Владимир Петрович}
\reviewerCredentials{к.т.н., доцент, ведущий научный сотрудник ИПМ им. М.В. Келдыша РАН.}

\headOfDepartmentTitle{И.о. директора Образовательного\\центра Института № 8}
\headOfDepartment{Крылов Сергей Сергеевич}

\headOfDepartmentTitle{Заведующий кафедрой 806}

\addbibresource{main.bib}

\makeatletter
\renewcommand{\makeSupervisorLine}{
    \vspace{-0.25cm}

    \textbf{Руководитель}\ \uline{Судаков Владимир Анатольевич, доктор технических наук, профессор, \hfill}\\
    \uline{профессор кафедры 806 \hfill}
}

\renewcommand{\makeHeader}{
    \noindent\begin{minipage}{0.05\textwidth}
        \includegraphics[scale=0.4]{inc/mai.pdf}
    \end{minipage}
    \hfill
    \begin{minipage}{0.85\textwidth}\raggedleft
        \begin{center}
            \fontsize{11pt}{10pt}\selectfont
            \textbf{МИНИСТЕРСТВО НАУКИ И ВЫСШЕГО ОБРАЗОВАНИЯ \\
                РОССИЙСКОЙ ФЕДЕРАЦИИ \\
                \hfill\break
                ФЕДЕРАЛЬНОЕ ГОСУДАРСТВЕННОЕ БЮДЖЕТНОЕ \\
                ОБРАЗОВАТЕЛЬНОЕ УЧРЕЖДЕНИЕ ВЫСШЕГО ОБРАЗОВАНИЯ \\
                <<МОСКОВСКИЙ АВИАЦИОННЫЙ ИНСТИТУТ
            } \\
            (национальный исследовательский университет)>>
        \end{center}
    \end{minipage}\\
    \noindent\parbox[t][24pt][c]{\textwidth}{\rule{\textwidth}{1.5pt}}
}
\makeatother

\begin{document}
\fillTaskHeader

\vspace{-0.6cm}

\SetSinglespace{1.2}
\singlespacing

\textbf{1.} \makeThemeLine

\textbf{2. Срок сдачи обучающимся законченной работы} \uline{24.05.2026\hfill}

\textbf{3. Задание и исходные данные к работе}\\
Разработать методы локального разрешения конфликтов для централизованного управления группой
автономных доставочных агентов, движущихся по графу дорожной сети с узкими местами, и
сопоставить их с классическими алгоритмами многоагентного планирования в задаче предотвращения
взаимоблокировок. Предусмотреть формальную постановку задачи, набор интегральных метрик
и симулятор городской среды. Исходными данными служат графы
дорожной сети, построенные по фрагментам реальной карты Москвы, и конфигурации агентов различной
плотности.

\textbf{Перечень иллюстративно-графических материалов:}
\vspace{\parskip}
\begin{table}
    \begin{tabularx}{\textwidth}{|c|X|c|}\hline
    \textbf{№ п/п} & \centering\textbf{Наименование} & \textbf{Количество листов} \\ \hline
    1 & Раздаточный материал &  \\ \hline
    \end{tabularx}
\end{table}

% Все что выше должно уместиться на одну страницу. В крайнем случая можно переносить части вниз (или наоборот)
\vfill\clearpage
\newgeometry{left=15mm, right=30mm, top=20mm, bottom=20mm} % зеркальные поля на второй странице

\textbf{4. Перечень подлежащих разработке разделов и этапы выполнения работы}
\vspace{\parskip}
\begin{table}
    \fontsize{10pt}{12pt}\selectfont
    \begin{tabularx}{\textwidth}{|c|>{\centering\arraybackslash}X|c|c|c|}\hline
    \textbf{\makecell{№\\п/п}} & \textbf{Наименование раздела или этапа} &
    \textbf{\makecell{Трудоёмкость в\\\% от полной\\трудоёмкости\\ВКР}} &
    \textbf{Срок выполнения} & \textbf{Примечание} \\ \hline
    1 & Изучение литературы и обзор существующих методов многоагентного планирования & 15\,\% & 02.09.2024--31.01.2025 & \\ \hline
    2 & Формальная постановка задачи и разработка набора интегральных метрик качества & 10\,\% & 01.02.2025--31.05.2025 & \\ \hline
    3 & Разработка методов локального разрешения конфликтов и адаптивной модификации A* & 25\,\% & 01.06.2025--31.10.2025 & \\ \hline
    4 & Разработка модульного симулятора и подготовка тестовых графов & 25\,\% & 01.11.2025--31.01.2026 & \\ \hline
    5 & Проведение сравнительного эксперимента и статистический анализ результатов & 15\,\% & 01.02.2026--31.03.2026 & \\ \hline
    6 & Оформление выпускной квалификационной работы & 10\,\% & 01.04.2026--25.05.2026 & \\ \hline
    \end{tabularx}
\end{table}

\textbf{5. Исходные материалы и пособия}\\ % основные работы, 25-30% от приведенного в работе списка
\fullcite{mapf-definition}\\
\fullcite{cbs}\\
\fullcite{ccbs}\\
\fullcite{sipp}\\
\fullcite{whca}\\
\fullcite{online-shortest-path}

\textbf{6. Дата выдачи задания} \uline{02.09.2024\hfill}

\makeTaskSignatures
\end{document}
